\input{../preamble.tex}

\SetDate[25/03/2026]

\reversemarginpar
\setlength{\marginparsep}{.5cm}

\begin{document}
\pagestyle{fancy}
\fancyhead[L]{Tle STMG}
\fancyhead[C]{\textbf{Évaluation blanche — Statistiques à deux variables}}
\fancyhead[R]{\today}

%\null\vspace{-30pt}
Consignes particulières : 
\begin{itemize}[label=$\bullet$]
	\item 
	La calculatrice est {autorisée}.
	\item
	L'évaluation fait \pageref{lastpage} page. La somme des points est \total{points}.
	\item
	Formules du coefficient directeur $a$ et de l'ordonnée à l'origine $b$ :
	\begin{align*}
		a = \dfrac{y_B-y_A}{x_B-x_A} && \et && b =y_A - ax_A.
	\end{align*}
\end{itemize}

\marginpar{[pts]}
\hrule



\exemulticols{6}{
	Deux ajustements affines sont proposés pour approximer le nuage de points ci-contre :
		\begin{align*}
			f(x) = 9-0,9x, && \et && g(x) = 9,5 - 1,1x.
		\end{align*}
	\begin{enumerate}
		\item
		Remplir les deux premières lignes du tableau, chaque colonne correspondant à un point du nuage de coordonnées $(x ; y)$.
		\item
		À quelle droite, $(d_1)$ ou $(d_2)$, correspond la fonction $f$ ? Justifier.
		\item
		Calculer l'erreur commise par chaque ajustement affine proposé à l'aide du tableau.
		\item
		Quel ajustement permet de mieux décrire les données ?
	\end{enumerate}
	\vfill
}{
	\null\vspace{-40pt}
	\begin{center}
	\begin{tikzpicture}
		\begin{axis}[xmin = 0, ymin=0, ymax=10, axis line style=->, extra y ticks={0}, extra x ticks={0}]
			\addplot[only marks, BLACK, mark size=2pt]  coordinates {(0,8)(1,9)(2,7)(3,6)(4,6)(6,2)(7,3)}; 
			
			\addplot[very thick, dashed, BLUE_E] expression[domain=0:7, samples=2]{9-.9*x} node[right]{$(d_2)$};
			\addplot[very thick, dashed, RED_E] expression[domain=0:7, samples=2]{9.5-1.1*x} node[right]{$(d_1)$};
		\end{axis}
	\end{tikzpicture}
	\end{center}
}{exe:l2-comparison}{
	\begin{enumerate}
		\item
		On note les coordonnées $(x;y)$ de chaque point du nuage dans le tableau ci-dessous.
		\item
		Comme $f(0) = 9$ et $g(0) = 9,5$, la droite $(d_2)$ correspond à la fonction $f$ car elle passe par $(0;9)$.
		\item
		On remplit le tableau à l'aide des expressions algébriques de $f$ et de $g$.
		\item
		On calcule la somme des erreurs individuelles pour obtenir l'erreur d'approximation associée à chaque fonction. 
		
		L'erreur d'approximation de $f$ est donc
			\[ 1 + 0,81 + 0,04 + 0,09 + 0,35 + 2,56 + 0,09 = 4,94. \]
		L'erreur d'approximation de $g$ est
			\[ 2,25 + 0,36 + 0,09 + 0,04 + 0,81 + 0,81 + 1,44 = 5,8. \]
		L'erreur la plus basse gagne : ici la fonction $f$ approxime mieux le nuage de point que la fonction $g$.
	\end{enumerate}
	
	\begin{center}
\def\arraystretch{1.3}
\setlength\tabcolsep{20pt}
	\begin{tabular}{|*{8}{c|}}\hline
		$x$ & 0 & 1 & 2 & 3 & 4 & 6 & 7
 		\\\hline
		$y$ & 8 & 9 & 7 & 6 & 6 & 2 & 3
		\\\hline
		$\bigl[f(x)-y\bigr]^2$ & 1 & 0,81 & 0,04 & 0,09 & 0,35 & 2,56 & 0,09
		\\\hline
		$\bigl[g(x)-y\bigr]^2$ & 2,25 & 0,36 & 0,09 & 0,04 & 0,81 & 0,81 & 1,44
		\\\hline
	\end{tabular}
	\end{center}
}
	\begin{center}
	
\def\arraystretch{1.3}
\setlength\tabcolsep{25pt}
	\begin{tabular}{|*{8}{c|}}\hline
		$x$ & & & & & & &
		\\\hline
		$y$ & & & & & & &
		\\\hline
		$\bigl[f(x)-y\bigr]^2$ & & & & & & &
		\\\hline
		$\bigl[g(x)-y\bigr]^2$ & & & & & & &
		\\\hline
	\end{tabular}
	\end{center}

\exe{6}{
	Pour chaque jeu de données ci-dessous, un ajustement affine est-il pertinent ? Justifier.
	
\newcommand{\scale}{.7}
	\begin{multicols}{3}
	\begin{center}
	\begin{tikzpicture}[scale=\scale]
		\begin{axis}[xmin = 0, xmax=13, ymin=1200, ymax=3500, ticks = none, clip=false, axis line style=->,]
			\addplot[only marks, ORANGE, mark size=2pt] table[x=anc, y=sal] {data/1-1.dat};
		\end{axis}
	\end{tikzpicture}
	
	(a)
	\end{center}
	
	\columnbreak
	
	\begin{center}
	\begin{tikzpicture}[scale=\scale]
		\begin{axis}[xmin = 0, xmax=13, ymin=1000, ymax=1500, ticks = none, clip=false, axis line style=->,]
			\addplot[only marks, ORANGE, mark size=2pt] table[x=anc, y=sal] {data/1-2.dat};
		\end{axis}
	\end{tikzpicture}
	
	(b)
	\end{center}
	
	\columnbreak
	
	\begin{center}
	\begin{tikzpicture}[scale=\scale]
		\begin{axis}[xmin =2, xmax=6, ymin=0, ymax=100, ticks = none, axis line style=->,]
			\addplot[only marks, ORANGE, mark size=2pt] table[x=anc, y=sal] {data/1-3.dat};
			%\addplot[very thick, dashed, BLUE_E] expression[domain=0:5, samples=2]{950 + 150*x};
			%\addplot[very thick, dashed, RED_E] expression[domain=0:5, samples=2]{2050 - 150*x};
		\end{axis}
	\end{tikzpicture}
	
	(c)
	\end{center}
	\end{multicols}
}{exe:pertinentYN}{
	\underline{Nuage (a)}
	
	Un ajustement affine n'est pas pertinent car les points ne suivent pas une tendance affine : ils ne suivent pas une droite.
	Cependant, la courbe tracée par le nuage ressemble à $y=x^2$.
	Une transformation de données semble donc pertinent ici : soit $x' = x^2$, soit $y' = \sqrt{y}$.
	Après tranformation, le nuage présentera peut-être une tendance affine.
	
	\underline{Nuage (b)}
	
	Oui, un ajustement affine est pertinent : les points sont plus ou moins alignés, avec tendance décroissante.
	
	\underline{Nuage (c)}
	
	Plusieurs ajustements affines semblent possibles ici, avec des erreurs d'approximations similaires.
	Cependant, les ajustements ont des tendances contradictoires, ce qui sous-entend qu'un ajustement affine n'est pas pertinent : il serait impossible d'extrapoler à l'aide de celui-ci.
	Le nuage présente un nombre insuffisant de point pour exhiber une tendance claire.
}

\newpage

\exe{8}{
	Leonardo Fibonacci élève des lapins de la race « Romarin » et étudie l'évolution de la population.
	Partant de l'année 0, il recense chaque année le nombre de lapins et le note dans son carnet.
	
	Ne sachant relire son écriture, les données du tableau sont malheureusement incomplètes.
	 
	\begin{multicols}{2}
	
	\begin{center}
	\begin{tikzpicture}
	\begin{axis}[
 		axis line style=->,
		grid=both,
		minor y tick num=1,
		xtick distance=1,
		ytick distance=30,
		extra x ticks={0},
		extra y ticks={0},
		x=20pt,
		y=.8pt,
		xmin=0,
		xmax=12,
		ymin=0,
		ymax=240,
		xlabel = Année,
		ylabel = Nombre de lapins,
		ylabel near ticks,
		xlabel near ticks,
		clip=true,
		]
		\addplot[only marks, transparent, mark size=2pt] table[x=n, y=Fn] {data/1-4.dat};
	\end{axis}
	\end{tikzpicture}
	\end{center}
	
\def\arraystretch{1.25}
\setlength\tabcolsep{5pt}
	\hspace{50pt}
	\begin{tabular}{|c|c|c|}
		\hline
		\textbf{Année} & \textbf{Lapins} & \textbf{LogLapins}
		\\\hline
		0 & 1 & \\ \hline
		1 & 1 & \\ \hline
		2 & 2 & \\\hline
		3 & 3 & \\\hline
		4 & 5 & \\\hline
		5 & 8 & \\\hline
		%6 & 13 & \\ \hline
		6 &  & \\ \hline
		7 & 21 & \\\hline
		%8 & 34 & \\\hline
		%9 & 55 & \\\hline
		8 &  & \\\hline
		9 &  & \\\hline
		10 & 89 & \\\hline
		11 & 144 & \\\hline
		12 & 233 & \\\hline
	\end{tabular}
	\vfill\null
	\end{multicols}
	
	\begin{enumerate}
		\item
		Représenter les données de Fibonacci dans le repère ci-dessus.
		\item
		Un ajustement affine des données vous semble-t-il pertinent ? Justifier.
		\item
		Fibonacci, suspectant une croissance exponentielle, décide d'appliquer le logarithme au nombre de lapins.
		Compléter le tableau arrondissant les logarithmes à $10^{-1}$ près.
		\item
		Représenter le nouveau nuage de points dans le repère ci-dessous.
		\item 
		Déterminer l'équation d'une droite $(d)$ d'ajustement affine liant le logarithme du nombre de lapins $\log(y)$ et l'année $x$ sous la forme $\log(y)=ax+b$.
		Arrondir les paramètres $a$ et $b$ à l'unité.
		\item
		Remplir les données manquantes du tableau de Fibonacci grâce à votre ajustement affine.
		\item
		Quelle sera la taille de la population de lapins à l'année 20 ?
		Prédire à l'aide de votre ajustement affine.
	\end{enumerate}
	
	
	\begin{center}
	\begin{tikzpicture}
	\begin{axis}[
 		axis line style=->,
		grid=both,
		minor y tick num=1,
		xtick distance=1,
		ytick distance=.5,
		extra x ticks={0},
		extra y ticks={0},
		x=20pt,
		%y=5pt,
		xmin=0,
		xmax=20,
		ymin=0,
		ymax=4,
		xlabel = Année,
		ylabel = Logarithme du nombre de lapins,
		ylabel near ticks,
		xlabel near ticks,
		clip=true,
		]
		\addplot[only marks, transparent, mark size=2pt] table[x=n, y=log] {data/1-4.dat};
		
		\addplot[thick, dashed, transparent] expression[domain=0:20, samples=2]{log10(1/2+sqrt(5)/2)*(x-.5)};
	\end{axis}
	\end{tikzpicture}
	\end{center}
	
	
	
}{exe:LuLu}{

	
	\begin{enumerate}
		\item\,
		
		\begin{center}
		\begin{tikzpicture}
		\begin{axis}[
	 		axis line style=->,
			grid=both,
			minor y tick num=1,
			xtick distance=1,
			ytick distance=30,
			extra x ticks={0},
			extra y ticks={0},
			x=20pt,
			y=.8pt,
			xmin=0,
			xmax=12,
			ymin=0,
			ymax=240,
			xlabel = Année,
			ylabel = Nombre de lapins,
			ylabel near ticks,
			xlabel near ticks,
			clip=true,
			]
			\addplot[only marks, ORANGE, mark size=2pt] table[x=n, y=Fn] {data/1-402.dat};
		\end{axis}
		\end{tikzpicture}
		\end{center}
		\item
		L'ajustement affine ne semble pas pertinent ici : les points ne sont pas (plus ou moins) alignés.
		\item\,
		
		\begin{center}
\def\arraystretch{1.25}
\setlength\tabcolsep{5pt}
		\begin{tabular}{|c|c|c|}
			\hline
			\textbf{Année} & \textbf{Lapins} & \textbf{LogLapins}
			\\\hline
			0 & 1 & 0 \\ \hline
			1 & 1 & 0 \\ \hline
			2 & 2 & 0,7 \\\hline
			3 & 3 & 1,1 \\\hline
			4 & 5 & 1,6 \\\hline
			5 & 8 & 2,1 \\\hline
			%6 & 13 & \\ \hline
			6 &  & \\ \hline
			7 & 21 & 3 \\\hline
			%8 & 34 & \\\hline
			%9 & 55 & \\\hline
			8 &  & \\\hline
			9 &  & \\\hline
			10 & 89 & 4,5 \\\hline
			11 & 144 & 5 \\\hline
			12 & 233 & 5,4 \\\hline
		\end{tabular}
		\end{center}
		\item\,
		
		\begin{center}
		\begin{tikzpicture}
		\begin{axis}[
	 		axis line style=->,
			grid=both,
			minor y tick num=1,
			xtick distance=1,
			ytick distance=.5,
			extra x ticks={0},
			extra y ticks={0},
			x=20pt,
			%y=5pt,
			xmin=0,
			xmax=20,
			ymin=0,
			ymax=4,
			xlabel = Année,
			ylabel = Logarithme du nombre de lapins,
			ylabel near ticks,
			xlabel near ticks,
			clip=true,
			]
			\addplot[only marks, ORANGE, mark size=2pt] table[x=n, y=log] {data/1-402.dat};
			
			\addplot[thick, dashed, BLUE_E] expression[domain=0:20, samples=2]{log10(1/2+sqrt(5)/2)*(x-.5)};
		\end{axis}
		\end{tikzpicture}
		\end{center}
		\item 
		En prenant deux points $A(x_A;y_A)$ et $B(x_B ; y_B)$ de la droite et en utilisant les formules données dans les consignes
		\begin{align*}
			a = \dfrac{y_B-y_A}{x_B-x_A} && \et && b =y_A - ax_A,
		\end{align*}
		on obtient environ
		\begin{align*}
			a = 0,2 && \et && b =-0,1,
		\end{align*}
		desquels nous déduisons que
			\[ \log(y) = 0,2x - 0,1. \]
		Il s'avère que la valeur exacte de $a$ est donnée par
			\[ a = \log\left(\dfrac{1+\sqrt5}2\right), \]
		le logarithme du nombre d'or ! (très fort)
		\item
		Il manque les données lorsque l'année $x$ vaut $6 ; 8 ; 9$. On remplace par ces valeurs pour obtenir les logarithmes du nombre de lapins suivants.
		\begin{align*}
			0,2\times6 - 0,1 &= 1,1 \\
			0,2\times8 - 0,1 &= 1,5 \\
			0,2\times9 - 0,1 &= 1,7
		\end{align*}
		En connaissant le logarithme du nombre de lapins, il suffit d'appliquer la fonction inverse pour obtenir le nombre de lapins.
		D'après le cours, $10^{\log(x)} = x$, donc le nombre de lapins aux années $6 ; 8 ; 9$ sont approximativement donnés par
		\begin{align*}
			10^{1,1} &\approx 12 \\
			10^{1,5} &\approx 32 \\
			10^{1,7} &\approx 50 \\
		\end{align*}
		\item
		À l'année 20, on pose $x=20$ pour obtenir
			\[ \log(y) = 0,2\times20-0,1 = 3,9. \]
		D'où
			\[ y = 10^{3,9} \approx 7~943. \]
		Nous pouvons donc prédire 7 943 lapins au bout de 20 ans !
	\end{enumerate}
		
		
}



%%%%%%%%%%%%

\label{lastpage}
\newpage
\fancyhead[C]{\textbf{Solutions}}
\shipoutAnswer

\end{document}
